\section{Complete Computations: Conditional Extensions}
\label{sec:examples_complete_conditional}

The following material extends the complete computations of
Section~\ref{sec:examples_complete} to the nonabelian setting.
The physics-bridge verification
(Proposition~\ref{prop:su2-hypotheses}) confirms that
$SU(2)$ Chern--Simons lies in the scope of
Theorem~\ref{thm:physics-bridge}; consequently, the results
below are now proved unconditionally.  The explicit
$\Ainf$ computation through arity~$4$, including the Jacobi
mechanism for $m_4 = 0$ on cohomology and tower termination at
$r_{\max} = 3$, appears in
Propositions~\ref{prop:su2-m3}--\ref{prop:su2-truncation}.

%%%%%%%%%%%%%%%%%%%%%%%%%%%%%%%%%%%%%%%%%%%%%%%%%%%%%%%%%%%%%%%%%%%%%%%%%%%%%
\subsection{Example IV: Nonabelian Chern--Simons ($SU(2)$ at level $k$)}
\label{ex:nonabelian_CS_complete}
\index{Chern--Simons!nonabelian|textbf}
\index{Kac--Moody algebra!affine sl2@affine $\mathfrak{sl}_2$}
%%%%%%%%%%%%%%%%%%%%%%%%%%%%%%%%%%%%%%%%%%%%%%%%%%%%%%%%%%%%%%%%%%%%%%%%%%%%%

The cubic self-coupling of the gauge field produces nonvanishing
$m_3$, and in the standard affine HT gauge realization the
open-colour Koszul dual is the Yangian $Y(\mathfrak{sl}_2)$.

\subsubsection{HT twist and field content}

Begin with 3d $SU(2)$ Chern--Simons theory on $\R_t \times \C_z$
at level~$k$.  The action is
\begin{equation}\label{eq:su2-cs-action}
S = \frac{k}{4\pi} \int_{\R \times \C}
  \operatorname{tr}\!\Bigl(
    A \wedge dA + \tfrac{2}{3}\, A \wedge A \wedge A
  \Bigr).
\end{equation}
After the HT twist and BV gauge fixing, the field content on
$\R_t \times \C_z$ is:
\begin{itemize}
\item Gauge field components
  $A_z^a(t,z)$ ($a = 1,2,3$ for
  $\mathfrak{sl}_2 = \operatorname{span}\{e,f,h\}$),
  of cohomological degree~$0$;
\item Antifield/ghost $c^a(t,z)$ of degree~$1$;
\item Lagrange multiplier
  $A_{\bar z}^a(t,z)$ of degree~$0$.
\end{itemize}

\subsubsection{BV-BRST differential and propagator}

The BV-BRST differential acts by
\begin{equation}\label{eq:su2-brst}
Q \cdot A_z^a = d_t c^a + f^a{}_{bc}\, A_z^b\, c^c,
\qquad
Q \cdot c^a = \tfrac{1}{2}\, f^a{}_{bc}\, c^b\, c^c,
\end{equation}
where $f^a{}_{bc}$ are the structure constants of
$\mathfrak{sl}_2$ (with $[e,f] = h$, $[h,e] = 2e$,
$[h,f] = -2f$).  The key difference from the abelian case:
$Q$ now has a \emph{quadratic} term from the gauge
self-interaction.

The HT propagator is the same as for abelian CS:
\begin{equation}\label{eq:su2-propagator}
\langle A_z^a(t_1,z_1)\, A_{\bar z}^b(t_2,z_2) \rangle
= \frac{\delta^{ab}}{2\pi} \cdot
  \frac{\Theta(t_1 - t_2)}{z_1 - z_2}.
\end{equation}
The Killing form $\delta^{ab}$ (normalized to
$\operatorname{tr}(t^a t^b) = \frac{1}{2}\delta^{ab}$
in the fundamental representation) couples the color indices.

\subsubsection{Binary operation $m_2$ and $\lambda$-bracket}

The binary operation $m_2$ is computed from the propagator
\eqref{eq:su2-propagator} via the same FM integral as in the
abelian case, with the OPE now producing color-dependent
residues.

\begin{proposition}[$\mathfrak{sl}_2$ $\lambda$-bracket;
\ClaimStatusProvedHere]
\label{prop:su2-lambda-bracket}
Setting
$J^a(z) := A_z^a(t,z)|_{t=0}$ (boundary current), the
$\lambda$-bracket is:
\begin{equation}\label{eq:su2-lambda}
\{J^a{}_\lambda J^b\}
= f^{ab}{}_c\, J^c + k\, \delta^{ab}\, \lambda.
\end{equation}
This is the defining relation of the affine Kac--Moody algebra
$\widehat{\mathfrak{sl}}_2$ at level~$k$.
\end{proposition}

\begin{proof}
The singular part of $m_2(J^a, J^b)$ is extracted from the
simple-pole residue of the propagator
$\langle A_z^a(z_1)\, A_{\bar z}^b(z_2)\rangle$ at $z_1 = z_2$.
The structure constant term $f^{ab}{}_c\, J^c$ comes from the
cubic vertex $f^a{}_{bc}\, A_z^b\, c^c$ in the BRST differential
\eqref{eq:su2-brst}: the tree-level Feynman diagram with one
internal vertex and two external legs produces a simple-pole
contribution proportional to $f^{ab}{}_c$.  The central term
$k\, \delta^{ab}\, \lambda$ comes from the double-pole residue of
the free propagator, exactly as in the abelian case.
\end{proof}

\subsubsection{The cubic operation $m_3$}

Unlike the abelian case, the nonabelian cubic vertex generates a
nonvanishing $m_3$.

\begin{proposition}[Nonabelian $m_3$: explicit formula; \ClaimStatusProvedHere]
\label{prop:su2-m3}
The cubic $\Ainf$ operation $m_3$ for
$\widehat{\mathfrak{sl}}_2$ at level~$k$ is nonzero at chain
level.  In the Chevalley basis $\{e, f, h\}$, the general
structural form of the Stasheff associator is
\begin{equation}\label{eq:su2-m3}
\begin{split}
m_3(J^a, J^b, J^c;\, \lambda_{12}, \lambda_{23})
&= \bigl(f^{ab}{}_d\, f^{dc}{}_g
  - f^{bc}{}_d\, f^{ad}{}_g\bigr)\, J^g \\
&\quad + \bigl(f^{ab}{}_d\, \delta^{dc}\,(\lambda_{12}+\lambda_{23})
  - f^{bc}{}_d\, \delta^{ad}\, \lambda_{12}\bigr)\, k,
\end{split}
\end{equation}
where the field-dependent part is the Jacobi expression
$[[J^a, J^b], J^c] - [J^a, [J^b, J^c]]$ (the difference of
the two $m_2 \circ m_2$ compositions), and the spectral-parameter
terms couple the central extension to the structure constants.
The nonvanishing triples are precisely those spanning a path
through the root system:
\begin{alignat}{2}
m_3(e, f, h) &= -2h + 2k\lambda_{23}, &\qquad
m_3(h, e, f) &= 2h + 2k\lambda_{23}, \notag \\
m_3(e, h, f) &= -2k\lambda_{23}, &\qquad
m_3(f, e, h) &= 2h + 2k\lambda_{23}, \notag \\
m_3(h, f, e) &= -2h + 2k\lambda_{23}, &\qquad
m_3(f, h, e) &= 2k\lambda_{23}. \label{eq:su2-m3-explicit}
\end{alignat}
All triples with repeated root generators give $m_3 = 0$.
On cohomology, every nonvanishing $m_3$ is $Q$-exact by Stokes on
$\FM_3(\C)$ (the AOS corner cancellation of
Section~\ref{sec:FM_calculus}): the Lie conformal Jacobi identity
ensures the associator is a coboundary, not a cocycle.
\end{proposition}

\begin{proof}
The explicit components~\eqref{eq:su2-m3-explicit} are computed by
evaluating the Stasheff associator
$A_3(J^a, J^b, J^c;\, \lambda_{12}, \lambda_{23})
= m_2(m_2(J^a, J^b;\, \lambda_{12}),\, J^c;\,
\lambda_{12} + \lambda_{23})
- m_2(J^a,\, m_2(J^b, J^c;\, \lambda_{23});\, \lambda_{12})$
on all independent generator triples.  The computation uses only
the $\lambda$-bracket~\eqref{eq:su2-lambda}, conformal
sesquilinearity, and the vacuum rule $m_2(\mathbf{1}, a) = 0$.
The nonvanishing triples are precisely those where the Jacobi
expression $f^{ab}{}_d f^{dc}{}_g - f^{bc}{}_d f^{ad}{}_g$ is
nonzero.  Triples with repeated root generators give $A_3 = 0$
by two mechanisms: adjacent repetitions vanish
because $m_2(e,e) = m_2(f,f) = 0$; non-adjacent repetitions
such as $(e,f,e)$ vanish because the two compositions in the
Stasheff difference cancel; explicitly,
$f^{ef}{}_h\, f^{he}{}_e = f^{fe}{}_h\, f^{eh}{}_e = 2$, so
the Jacobi expression
$f^{ef}{}_h f^{he}{}_e - f^{fe}{}_h f^{eh}{}_e = 0$ and both
spectral terms vanish by $\delta^{he} = 0$.

The chain-level $m_3$ is obtained by antidifferentiation on
$\FM_3(\C) \cong [0,1]$; the specific
$\lambda_{23}$ dependence in~\eqref{eq:su2-m3-explicit}
corresponds to the left-endpoint gauge.  On cohomology,
the Stokes argument on $\FM_3(\C)$ shows each $m_3$ is
$Q$-exact: the boundary decomposes into two codimension-$1$
faces ($z_1 \to z_2$ and $z_2 \to z_3$), whose contributions
are the compositions $m_2 \circ m_2$.  The Lie conformal Jacobi
identity ensures $A_3 = d(m_3)$ with $m_3$ in the image of $d$,
so $[m_3] = 0$ in $H^\bullet$.
\end{proof}

\subsubsection{Higher operations $m_{k \geq 4}$}

\begin{proposition}[$\Ainf$ tower termination at arity~$3$;
\ClaimStatusProvedHere]
\label{prop:su2-truncation}
For $\widehat{\mathfrak{sl}}_2$ at level~$k$:
\begin{enumerate}[label=\textup{(\roman*)}]
\item $m_4 = 0$ on cohomology: the quartic Stasheff obstruction
  $A_4(J^a, J^b, J^c, J^d)$ vanishes identically by the Jacobi
  identity applied twice.  The field-dependent part assembles
  iterated structure constant contractions
  $f^{ab}{}_p f^{pc}{}_q f^{qd}{}_r$ whose Stasheff-signed sum
  is the generalised Jacobi identity at depth~$3$; the central
  terms cancel by the residue theorem on $\FM_4(\C)$.
\item $m_4 \neq 0$ at chain level: the square Feynman diagram
  with three cubic vertices contributes degree~$2$, matching
  $\FM_4(\C)$.  This chain-level contribution is $Q$-exact.
\item $m_{k \geq 5} = 0$ at all levels by AOS degree counting:
  a tree with $k$ external legs requires $k-2$ vertices of
  degree~$1$ each; for $k \geq 5$, the logarithmic form degree
  exceeds the dimension of $\FM_k(\C)$.
\item On cohomology $H^\bullet(\cA, Q)$, all $m_{k \geq 3}$
  vanish ($\Ainf$-formality): the PVA on $H^\bullet$ is
  $\widehat{\mathfrak{sl}}_2$ at level~$k$ with the standard
  $\lambda$-bracket.  Shadow depth $r_{\max} = 3$
  (class~$\mathbf{L}$).
\end{enumerate}
\end{proposition}

\begin{proof}
Part~(i): the quartic Stasheff obstruction on the quadruple
$(e, f, e, f)$ (passing through all three root spaces) is computed
term by term.  Of the five Stasheff terms, two vanish because
$m_3(e,f,e) = m_3(f,e,f) = 0$ (repeated root spaces); the
remaining three produce field-dependent parts $+2h$ and $-2h$ that
cancel, plus spectral-parameter terms that cancel by the residue
theorem.  For general quadruples, the field-dependent part
assembles into the Stasheff-signed sum of iterated contractions
$f^{ab}{}_p f^{pc}{}_q f^{qd}{}_r$, which vanishes by applying
the Jacobi identity $f^{ab}{}_p f^{pc}{}_q = f^{ac}{}_p f^{bp}{}_q
+ f^{bc}{}_p f^{ap}{}_q$ and then pairwise cancellation.

Part~(ii): the diagram with $k=4$ external legs and $2$ cubic
vertices has degree $2$, matching $\FM_4(\C)$.  The structure
constant contraction $f^{a}{}_{de}\, f^{e}{}_{fg}\, f^{fbc}$
is generically nonvanishing.

Part~(iii): a tree-level diagram with $k$ legs and $v = k - 2$
trivalent internal vertices has total degree $k-2$.  For
$k \geq 5$, the AOS relations force the integral to vanish.

Part~(iv): by induction on arity.  Base case: $m_4^H = 0$ from
part~(i).  Inductive step: for $k \geq 5$, the Stasheff relation
involves only $m_2$ and $m_3$ (all higher operations vanish by
hypothesis).  The compositions $m_3 \circ m_3$ at arity $5$
produce iterated structure constant contractions of depth~$4$,
which vanish by the same Jacobi mechanism.
\end{proof}

\subsubsection{Cohomology and PVA structure}

\begin{theorem}[PVA for nonabelian CS; \ClaimStatusProvedHere]
\label{thm:su2-pva}
Let\/ $\cA$ be a logarithmic\/ $\SCchtop$-algebra.  The cohomology
$H^\bullet(\A_{\text{CS}}, Q)$ is the affine Kac--Moody PVA
$\widehat{\mathfrak{sl}}_2$ at level~$k$: a
$(-1)$-shifted Poisson vertex algebra with
\begin{enumerate}[label=\textup{(\roman*)}]
\item generators $J^a$ ($a = 1,2,3$) of degree~$0$;
\item commutative product (regular part of $m_2$);
\item $\lambda$-bracket \eqref{eq:su2-lambda}.
\end{enumerate}
The PVA axioms (sesquilinearity, skewsymmetry, Jacobi, Leibniz)
hold by the general descent theorem
(Theorem~\ref{thm:cohomology_PVA}).
\end{theorem}

\subsubsection{Koszul dual and Yangian structure}

Assuming the affine example lies on the chirally Koszul locus,
Theorem~\ref{thm:lines_as_modules} identifies
$\mathcal{C}_{\mathrm{line}} \simeq \cA^!_{\mathrm{line}}\text{-}\mathbf{mod}$,
so the open-colour Koszul dual $\cA^!_{\mathrm{line}}$ controls line
operators.

\begin{theorem}[Koszul dual is $Y(\mathfrak{sl}_2)$;
\ClaimStatusProvedHere]
\label{thm:su2-yangian}
For the standard affine HT gauge realization with closed colour
$\widehat{\mathfrak{sl}}_2$ at level~$k$, the open-colour Koszul dual is the Yangian
$Y(\mathfrak{sl}_2)$, a dg-shifted Yangian
(Theorem~\ref{thm:Koszul_dual_Yangian}) with:
\begin{enumerate}[label=\textup{(\roman*)}]
\item RTT presentation: $R_{12}(u-v)\, T_1(u)\, T_2(v)
  = T_2(v)\, T_1(u)\, R_{12}(u-v)$;
\item spectral $R$-matrix from the Laplace transform
  (Proposition~\ref{prop:field-theory-r}):
  \begin{equation}\label{eq:su2-rmatrix}
  r(z) = \frac{k\,\Omega}{z},
  \qquad
  \Omega = \sum_{a=1}^3 t^a \otimes t^a
  \in \mathfrak{sl}_2 \otimes \mathfrak{sl}_2,
  \end{equation}
  where $\Omega$ is the Casimir element and $k$ is the affine level; at $k = 0$ the $r$-matrix vanishes identically (AP126 check);
\item the classical $r$-matrix \eqref{eq:su2-rmatrix} satisfies
  the classical Yang--Baxter equation
  $[r_{12}, r_{13}] + [r_{12}, r_{23}] + [r_{13}, r_{23}] = 0$.
\end{enumerate}
\end{theorem}

\begin{proof}
The $\lambda$-bracket \eqref{eq:su2-lambda} has singular part
$f^{ab}{}_c\, J^c / (z_1 - z_2) + k\, \delta^{ab} / (z_1 - z_2)^2$.
The Laplace transform
$r^{ab}(z) = \int_0^\infty e^{-\lambda z}\, \{J^a{}_\lambda J^b\}\, d\lambda$
gives
\[
r^{ab}(z) = \frac{f^{ab}{}_c\, J^c}{z} + \frac{k\, \delta^{ab}}{z^2}.
\]
The $k\,\delta^{ab}/z^2$ term is central (proportional to the Killing
form tensor) and drops out of the CYBE\@; the surviving term
$k\,\delta^{ab}/z^2$ gives, after $d\log$ absorption, the collision $r$-matrix $r(z) = k\,\Omega/z$, the standard rational
$r$-matrix for $\widehat{\mathfrak{sl}}_2$ at level~$k$.  The CYBE follows from the
Jacobi identity of $\mathfrak{sl}_2$ (the same argument as
Theorem~\ref{thm:YBE} restricted to the $3$-dimensional case).

The open-colour Yangian identification
$\cA^!_{\mathrm{line}} \simeq Y(\mathfrak{sl}_2)$ is the
affine specialization of
Theorem~\ref{thm:Koszul_dual_Yangian}; the RTT presentation is
the one encoded by the spectral $R$-matrix
\eqref{eq:su2-rmatrix}.
\end{proof}

\subsubsection{Genus-$1$ curvature}

\begin{proposition}[Curvature for $\widehat{\mathfrak{sl}}_2$;
\ClaimStatusProvedHere]
\label{prop:su2-curvature}
The genus-$1$ curvature
(Volume~I, Theorem~D, the modular characteristic) evaluates to
\begin{equation}\label{eq:su2-curvature}
\kappa(\widehat{\mathfrak{sl}}_2, k) = \frac{3(k + 2)}{4},
\end{equation}
where $3 = \dim(\mathfrak{sl}_2)$ and $h^\vee = 2$ is the
dual Coxeter number.  By complementarity (Volume~I, Theorem~C):
\begin{equation}
\kappa(\widehat{\mathfrak{sl}}_{2,k})
+ \kappa(\widehat{\mathfrak{sl}}_{2,-k-4}) = 0.
\end{equation}
The Koszul dual level is $k' = -k - 2h^\vee = -k - 4$
(Feigin--Frenkel involution).
\end{proposition}

\begin{proof}
The central charge of $\widehat{\mathfrak{sl}}_2$ at level~$k$
is $c = 3k/(k+2)$ (Sugawara formula with $\dim\mathfrak{sl}_2 = 3$
and $h^\vee = 2$).  The curvature $\kappa = \dim\mathfrak{g}\cdot(k+h^\vee)/(2h^\vee) = 3(k+2)/4$, matching Volume~I's computation.
Under $k \mapsto k' = -k - 2h^\vee = -k - 4$:
$\kappa(k') = 3(k'+2)/4 = 3(-k-4+2)/4 = 3(-k-2)/4 = -3(k+2)/4 = -\kappa(k)$,
confirming complementarity $\kappa + \kappa' = 0$.
\end{proof}

\subsubsection{Verification of the physics-bridge inputs}
\label{subsubsec:su2-verification}

\begin{proposition}[Physics-bridge verification for $SU(2)$ CS;
\ClaimStatusProvedHere]
\label{prop:su2-hypotheses}
$SU(2)$ Chern--Simons theory at integer level~$k \geq 1$ lies in the
scope of Theorem~\ref{thm:physics-bridge}.  Concretely:
\begin{enumerate}[label=\textup{(\roman*)}]
\item BV data and one-loop finiteness: the action \eqref{eq:su2-cs-action} is a
  classical BV theory with well-defined HT twist.  Gauge fixing
  produces a propagator with standard properties.  One-loop
  finiteness holds because 3d CS is topological: the only
  UV divergence is a one-loop shift $k \to k + h^\vee$
  (renormalization of the level), which is finite.
\item The propagator \eqref{eq:su2-propagator} is meromorphic in~$z$
  (simple pole at $z_1 = z_2$), exponential decay in~$t$
  (Heaviside step function).
\item The interaction vertex $f^a{}_{bc}$
  is polynomial in the fields, so Feynman integrands are
  logarithmic forms on $\FM_k(\C)$.  AOS relations hold by the
  general theory of configuration space integrals for
  polynomial interactions.  Stokes exactness holds on
  $\FM_k(\C)$ by the compactification's smooth-with-corners
  structure.
\end{enumerate}
Therefore Theorem~\ref{thm:physics-bridge} applies, and the resulting
observables form a logarithmic $\SCchtop$-algebra.  The induced
factorization compatibility is then part of the resulting
$C_\ast(W(\SCchtop))$-algebra structure.
\end{proposition}

\subsection{Summary Table Including Conditional Examples}

\begin{center}
\begin{tabular}{|l|c|c|c|c|}
\hline
\textbf{Theory} & $m_2$ & $m_3$ & $m_{k \geq 4}$ & $H^\bullet(\A, Q)$ \\
\hline
$SU(2)$ CS & $\frac{f^{ab}{}_c J^c}{z_{12}} + \frac{k\delta^{ab}}{z_{12}^2}$ & $\neq 0$ (explicit) & $m_4^H = 0$ (Jacobi); $m_{\geq 5} = 0$ & $\widehat{\mathfrak{sl}}_2$ at level $k$ \\
\hline
\end{tabular}
\end{center}
